\documentclass[12pt, a4paper]{article}

% ==========================================
% 1. COMPILER, LANGUAGE & FONT SETUP
% ==========================================
\usepackage{fontspec}
\usepackage{polyglossia}
\setmainlanguage{bengali}
\setotherlanguage{english}
\newfontfamily\bengalifont[Script=Bengali, AutoFakeBold=2.5, AutoFakeSlant=0.2]{Kalpurush}

% ==========================================
% 2. MATHEMATICS, UNITS & FORMATTING
% ==========================================
\usepackage{amsmath, amssymb, siunitx}
\usepackage[margin=1in]{geometry}
\usepackage{setspace}
\usepackage{enumitem}
\usepackage{microtype} % For better typography
\onehalfspacing % Better line spacing for readability

% ==========================================
% 3. COLORS & BEAUTIFICATION PACKAGES
% ==========================================
\usepackage[dvipsnames, table]{xcolor}
\usepackage[skins, breakable, theorems]{tcolorbox}
\usepackage{tikz}
\usetikzlibrary{shadows, arrows.meta, positioning, decorations.pathmorphing, calc}
\renewcommand{\labelitemi}{$\bullet$}

% Define custom beautiful colors
\definecolor{primary}{HTML}{0056b3}
\definecolor{secondary}{HTML}{e7f1ff}
\definecolor{success}{HTML}{198754}
\definecolor{successbg}{HTML}{d1e7dd}
\definecolor{accent}{HTML}{fd7e14}
\definecolor{darkgray}{HTML}{343a40}

% Custom Tcolorbox Styles
\tcbset{
    questionbox/.style={
        enhanced, colback=secondary, colframe=primary, arc=2mm, boxrule=1pt,
        fonttitle=\bfseries\Large, title=#1,
        drop fuzzy shadow=black!10,
        attach boxed title to top left={xshift=5mm, yshift=-3mm},
        boxed title style={colback=primary, arc=1mm}
    },
    answerbox/.style={
        enhanced, colback=successbg, colframe=success, arc=1.5mm, boxrule=1pt,
        left=2mm, right=2mm, top=2mm, bottom=2mm,
        drop fuzzy shadow=black!10
    },
    solutionbox/.style={
        enhanced, colback=white, colframe=darkgray, arc=2mm, boxrule=1pt,
        fonttitle=\bfseries\large, title=#1, breakable,
        coltitle=white, colbacktitle=darkgray,
        drop fuzzy shadow=black!10,
        attach boxed title to top center={yshift=-3mm},
        boxed title style={arc=1mm}
    }
}

\begin{document}

\newpage
% ------------------------------------------
% QUESTION 13 SECTION
% ------------------------------------------
\begin{tcolorbox}[questionbox={প্রশ্ন (Question) 1}]
    \noindent
    $10\text{ L}$ আয়তনের বদ্ধ পাত্রে $300\text{ K}$ তাপমাত্রায় $16\text{ g}$ অক্সিজেন গ্যাস চাপ প্রদর্শন করে। একই পাত্রে একই তাপমাত্রায় কত গ্রাম নাইট্রোজেন রাখলে একই চাপ প্রদর্শন করবে? (যাতে মোট চাপ দ্বিগুণ হয়)
    
    \vspace{0.8em}
    \begin{english}
    \noindent\textit{\color{darkgray}
    A closed vessel of volume $10\text{ L}$ contains $16\text{ g}$ of Oxygen gas at $300\text{ K}$ exerting a certain pressure. How many grams of Nitrogen gas must be placed in the same vessel at the same temperature to exert the same additional pressure (so that total pressure doubles)?
    }
    \end{english}
\end{tcolorbox}

\begin{itemize}[label={}, leftmargin=1cm, itemsep=0.5em]
    \item[\textbf{A)}] $14\text{ g}$
    \item[\textbf{B)}] $16\text{ g}$
    \item[\textbf{C)}] $18\text{ g}$
    \item[\textbf{D)}] $32\text{ g}$
\end{itemize}

\vspace{0.5em}
\begin{tcolorbox}[answerbox]
    \textbf{\color{success} উত্তর:} A) $14\text{ g}$
\end{tcolorbox}
\vspace{1em}

\begin{tcolorbox}[solutionbox={সমাধান (Solution)}]
    \noindent
    \textbf{ধাপ ১: অক্সিজেনের মোল সংখ্যা নির্ণয়} \\
    অক্সিজেনের ভর $m_1 = 16\text{ g}$ এবং আণবিক ভর $M_1 = 32\text{ g mol}^{-1}$।
    \[ n_1 = \frac{m_1}{M_1} = \frac{16}{32} = 0.5\text{ mol} \]
    
    \vspace{0.4em}
    \noindent\textbf{ধাপ ২: মূলনীতি প্রয়োগ} \\
    আদর্শ গ্যাস সমীকরণ $PV = nRT$ অনুযায়ী স্থির আয়তন ও তাপমাত্রায় চাপ মোল সংখ্যার সমানুপাতিক ($P \propto n$)। সমপরিমাণ অতিরিক্ত চাপ সৃষ্টি করতে হলে সমপরিমাণ মোল সংখ্যা প্রয়োজন। অর্থাৎ নাইট্রোজেনের মোল সংখ্যাও অক্সিজেনের সমান হতে হবে:
    \[ n_2 = n_1 = 0.5\text{ mol} \]
    
    \vspace{0.4em}
    \noindent\textbf{ধাপ ৩: প্রয়োজনীয় নাইট্রোজেনের ভর নির্ণয়} \\
    নাইট্রোজেনের আণবিক ভর $M_2 = 28\text{ g mol}^{-1}$।
    \begin{align*}
        m_2 &= n_2 \times M_2 \\
        &= 0.5 \times 28 \\
        &= \mathbf{14\text{ g}}
    \end{align*}
\end{tcolorbox}

\vspace{2em}
% ------------------------------------------
% QUESTION 14 SECTION
% ------------------------------------------
\begin{tcolorbox}[questionbox={প্রশ্ন (Question) 2}]
    \noindent
    $0^\circ\text{C}$ তাপমাত্রায় কোনো গ্যাসের চাপ $3 \times 10^5\text{ Pa}$। আয়তন স্থির রেখে $60^\circ\text{C}$ তাপমাত্রায় এর চাপ কত হবে?
    
    \vspace{0.8em}
    \begin{english}
    \noindent\textit{\color{darkgray}
    The pressure of a gas at $0^\circ\text{C}$ is $3 \times 10^5\text{ Pa}$. Keeping the volume constant, what will be its pressure at $60^\circ\text{C}$?
    }
    \end{english}
\end{tcolorbox}

\begin{itemize}[label={}, leftmargin=1cm, itemsep=0.5em]
    \item[\textbf{A)}] $4.66 \times 10^5\text{ Pa}$
    \item[\textbf{B)}] $3.66 \times 10^5\text{ Pa}$
    \item[\textbf{C)}] $2.66 \times 10^5\text{ Pa}$
    \item[\textbf{D)}] $5.67 \times 10^5\text{ Pa}$
\end{itemize}

\vspace{0.5em}
\begin{tcolorbox}[answerbox]
    \textbf{\color{success} উত্তর:} B) $3.66 \times 10^5\text{ Pa}$
\end{tcolorbox}
\vspace{1em}

\begin{tcolorbox}[solutionbox={সমাধান (Solution)}]
    \noindent
    \textbf{ধাপ ১: স্থির আয়তনে চাপের সূত্র} \\
    গে-লুসাকের চাপের সূত্র বা স্থির আয়তনের সূত্রানুসারে:
    \[ \frac{P_1}{T_1} = \frac{P_2}{T_2} \]
    
    \vspace{0.4em}
    \noindent\textbf{ধাপ ২: উপাত্তসমূহ এবং তাপমাত্রা রূপান্তর}
    \begin{itemize}[leftmargin=1.5em, itemsep=0.2em]
        \item আদি চাপ, $P_1 = 3 \times 10^5\text{ Pa}$
        \item আদি তাপমাত্রা, $T_1 = 273 + 0 = 273\text{ K}$
        \item অন্তিম তাপমাত্রা, $T_2 = 273 + 60 = 333\text{ K}$
    \end{itemize}
    
    \vspace{0.4em}
    \noindent\textbf{ধাপ ৩: অন্তিম চাপ $P_2$ হিসাব}
    \begin{align*}
        P_2 &= \frac{P_1}{T_1} \times T_2 \\
        &= \frac{3 \times 10^5}{273} \times 333 \\
        &= 1098.901 \times 333 \\
        &= \mathbf{3.66 \times 10^5\text{ Pa}}
    \end{align*}
\end{tcolorbox}

\vspace{2em}

% ------------------------------------------
% QUESTION 15 SECTION
% ------------------------------------------
\begin{tcolorbox}[questionbox={প্রশ্ন (Question) 3}]
    \noindent
    $20\text{ g}$ হিলিয়াম গ্যাস পূর্ণ একটি বেলুনের আয়তন $0.12\text{ m}^3$। বেলুনের ভেতরে গ্যাসের চাপ $1.5 \times 10^5\text{ N m}^{-2}$। বেলুনের ভেতরে গ্যাসের তাপমাত্রা কত?
    
    \vspace{0.8em}
    \begin{english}
    \noindent\textit{\color{darkgray}
    A balloon filled with $20\text{ g}$ of Helium gas has a volume of $0.12\text{ m}^3$. The gas pressure inside the balloon is $1.5 \times 10^5\text{ N m}^{-2}$. What is the temperature of the gas inside the balloon?
    }
    \end{english}
\end{tcolorbox}

\begin{itemize}[label={}, leftmargin=1cm, itemsep=0.5em]
    \item[\textbf{A)}] $300\text{ K}$
    \item[\textbf{B)}] $433.004\text{ K}$
    \item[\textbf{C)}] $273\text{ K}$
    \item[\textbf{D)}] $500\text{ K}$
\end{itemize}

\vspace{0.5em}
\begin{tcolorbox}[answerbox]
    \textbf{\color{success} উত্তর:} B) $433.004\text{ K}$
\end{tcolorbox}
\vspace{1em}

\begin{tcolorbox}[solutionbox={সমাধান (Solution)}]
    \noindent
    \textbf{ধাপ ১: আদর্শ গ্যাস সমীকরণ প্রয়োগ} \\
    আদর্শ গ্যাস সমীকরণটি হলো:
    \[ PV = nRT = \frac{m}{M} R T \]
    
    \vspace{0.4em}
    \noindent\textbf{ধাপ ২: প্রদত্ত উপাত্তসমূহ চিহ্নিতকরণ}
    \begin{itemize}[leftmargin=1.5em, itemsep=0.2em]
        \item গ্যাসের ভর, $m = 20\text{ g}$
        \item হিলিয়ামের আণবিক ভর, $M = 4\text{ g mol}^{-1}$
        \item মোল সংখ্যা, $n = \frac{20}{4} = 5\text{ mol}$
        \item আয়তন, $V = 0.12\text{ m}^3$
        \item চাপ, $P = 1.5 \times 10^5\text{ N m}^{-2}$
        \item মোলার গ্যাস ধ্রুবক, $R = 8.314\text{ J mol}^{-1}\text{ K}^{-1}$
    \end{itemize}
    
    \vspace{0.4em}
    \noindent\textbf{ধাপ ৩: তাপমাত্রা $T$ হিসাব}
    \begin{align*}
        T &= \frac{P V}{n R} \\
        &= \frac{1.5 \times 10^5 \times 0.12}{5 \times 8.314} \\
        &= \frac{18000}{41.57} \\
        &= \mathbf{433.004\text{ K}}
    \end{align*}
\end{tcolorbox}
% ------------------------------------------
% QUESTION 17 SECTION
% ------------------------------------------
\begin{tcolorbox}[questionbox={প্রশ্ন (Question) 4}]
    \noindent
    একটি এয়ার কন্ডিশনার $40^\circ\text{C}$ তাপমাত্রা এবং $75\%$ আপেক্ষিক আর্দ্রতা বিশিষ্ট বাইরের বায়ুকে টেনে ঘরের ভেতরে আনে এবং $20^\circ\text{C}$ তাপমাত্রা এবং $54\%$ আপেক্ষিক আর্দ্রতা বিশিষ্ট বাতাসে পরিণত করে। এই প্রক্রিয়ায় প্রতি ঘনমিটারে কতটা জলীয় বাষ্প ঘনীভূত হবে? ($40^\circ\text{C}$ এবং $20^\circ\text{C}$ তাপমাত্রায় সম্পৃক্ত জলীয় বাষ্পের ঘনত্ব যথাক্রমে $45\text{ g m}^{-3}$ এবং $17\text{ g m}^{-3}$)
    
    \vspace{0.8em}
    \begin{english}
    \noindent\textit{\color{darkgray}
    An air conditioner pulls outdoor air at a temperature of $40^\circ\text{C}$ and $75\%$ relative humidity into a room and converts it into air at $20^\circ\text{C}$ and $54\%$ relative humidity. How much water vapour will condense per cubic meter in this process? (Saturated vapour density at $40^\circ\text{C}$ and $20^\circ\text{C}$ are $45\text{ g m}^{-3}$ and $17\text{ g m}^{-3}$ respectively)
    }
    \end{english}
\end{tcolorbox}

\begin{itemize}[label={}, leftmargin=1cm, itemsep=0.5em]
    \item[\textbf{A)}] $33.75\text{ g m}^{-3}$
    \item[\textbf{B)}] $24.57\text{ g m}^{-3}$
    \item[\textbf{C)}] $9.18\text{ g m}^{-3}$
    \item[\textbf{D)}] $12.50\text{ g m}^{-3}$
\end{itemize}

\vspace{0.5em}
\begin{tcolorbox}[answerbox]
    \textbf{\color{success} উত্তর:} B) $24.57\text{ g m}^{-3}$
\end{tcolorbox}
\vspace{1em}

\begin{tcolorbox}[solutionbox={সমাধান (Solution)}]
    \noindent
    \textbf{ধাপ ১: সূত্র উপস্থাপন} \\
    জলীয় বাষ্পের প্রকৃত ঘনত্ব $f = \text{RH} \times F$ (যেখানে $\text{RH}$ হলো আপেক্ষিক আর্দ্রতা এবং $F$ হলো সম্পৃক্ত জলীয় বাষ্পের ঘনত্ব)।
    
    \vspace{0.4em}
    \noindent\textbf{ধাপ ২: প্রাথমিক ও চূড়ান্ত ঘনত্ব নির্ণয়}
    \begin{itemize}[leftmargin=1.5em, itemsep=0.2em]
        \item বাইরের বায়ুতে জলীয় বাষ্পের প্রাথমিক ঘনত্ব, $f_1 = 0.75 \times 45\text{ g m}^{-3} = 33.75\text{ g m}^{-3}$
        \item ঘরের ভেতরের বায়ুতে জলীয় বাষ্পের চূড়ান্ত ঘনত্ব, $f_2 = 0.54 \times 17\text{ g m}^{-3} = 9.18\text{ g m}^{-3}$
    \end{itemize}
    
    \vspace{0.4em}
    \noindent\textbf{ধাপ ৩: ঘনীভূত পানির পরিমাণ হিসাব}
    \begin{align*}
        \Delta f &= f_1 - f_2 \\
        &= 33.75 - 9.18 \\
        &= \mathbf{24.57\text{ g m}^{-3}}
    \end{align*}
\end{tcolorbox}

\vspace{2em}

% ------------------------------------------
% QUESTION 18 SECTION
% ------------------------------------------
\begin{tcolorbox}[questionbox={প্রশ্ন (Question) 5}]
    \noindent
    একটি পুকুরের পানির গভীরতা $6\text{ m}$। বায়ুমণ্ডলের তাপমাত্রা $27^\circ\text{C}$ এবং পানির ভেতরে নিমজ্জনের সাথে প্রতি মিটার গভীরতার জন্য তাপমাত্রা $0.5^\circ\text{C}$ হ্রাস পায়। পানির ঘনত্বের পরিবর্তন উপেক্ষা করলে, পুকুরের তলদেশে উৎপন্ন একটি বুদ্বুদ উপরিপৃষ্ঠে পৌঁছানোর সময় তার আয়তনের শতকরা পরিবর্তন কত হবে?
    
    \vspace{0.8em}
    \begin{english}
    \noindent\textit{\color{darkgray}
    The depth of a pond is $6\text{ m}$. The atmospheric temperature is $27^\circ\text{C}$ and inside the water, the temperature decreases by $0.5^\circ\text{C}$ per meter of depth. Neglecting changes in water density, what will be the percentage change in volume of a bubble produced at the bottom of the pond by the time it reaches the surface?
    }
    \end{english}
\end{tcolorbox}

\begin{itemize}[label={}, leftmargin=1cm, itemsep=0.5em]
    \item[\textbf{A)}] আয়তন $42.5\%$ বৃদ্ধি পাবে
    \item[\textbf{B)}] আয়তন $59.6\%$ বৃদ্ধি পাবে
    \item[\textbf{C)}] আয়তন $35.2\%$ বৃদ্ধি পাবে
    \item[\textbf{D)}] আয়তন $59.6\%$ হ্রাস পাবে
\end{itemize}

\vspace{0.5em}
\begin{tcolorbox}[answerbox]
    \textbf{\color{success} উত্তর:} B) আয়তন $59.6\%$ বৃদ্ধি পাবে
\end{tcolorbox}
\vspace{1em}

\begin{tcolorbox}[solutionbox={সমাধান (Solution)}]
    \noindent
    \textbf{ধাপ ১: উপরিপৃষ্ঠ ও তলদেশের অবস্থা নির্ধারণ}
    \begin{itemize}[leftmargin=1.5em, itemsep=0.2em]
        \item উপরিপৃষ্ঠের তাপমাত্রা, $T_2 = 273 + 27 = 300\text{ K}$
        \item উপরিপৃষ্ঠের চাপ, $P_2 = P_a = 1.013 \times 10^5\text{ N m}^{-2}$
        \item তলদেশের তাপমাত্রা, $T_1 = 273 + (27 - 6 \times 0.5) = 273 + 24 = 297\text{ K}$
        \item তলদেশের চাপ, $P_1 = P_a + h \rho g = 1.013 \times 10^5 + 6 \times 1000 \times 9.8 = 1.601 \times 10^5\text{ N m}^{-2}$
    \end{itemize}
    
    \vspace{0.4em}
    \noindent\textbf{ধাপ ২: আদর্শ গ্যাসের সমন্বয় সূত্র প্রয়োগ}
    \begin{align*}
        \frac{P_1 V_1}{T_1} &= \frac{P_2 V_2}{T_2} \\
        \implies \frac{V_2}{V_1} &= \frac{P_1 T_2}{P_2 T_1} \\
        &= \frac{1.601 \times 10^5 \times 300}{1.013 \times 10^5 \times 297} = 1.596
    \end{align*}
    
    \vspace{0.4em}
    \noindent\textbf{ধাপ ৩: শতকরা পরিবর্তন হিসাব}
    \begin{align*}
        \text{Percentage Change} &= \frac{V_2 - V_1}{V_1} \times 100\% \\
        &= (1.596 - 1) \times 100\% \\
        &= \mathbf{59.6\% \text{ (বৃদ্ধি)}}
    \end{align*}
\end{tcolorbox}
% ------------------------------------------
% QUESTION 34 SECTION
% ------------------------------------------
\begin{tcolorbox}[questionbox={প্রশ্ন (Question) 6}]
    \noindent
    কোনো স্থানে হাইগ্রোমিটারের শুষ্ক বাল্বের পাঠ $24^\circ\text{C}$ এবং শিশিরাঙ্ক $11.5^\circ\text{C}$। $24^\circ\text{C}$ তাপমাত্রায় গ্লাইশারের উৎপাদক $1.72$ হলে সিক্ত বাল্বের পাঠ কত? এবং $24^\circ\text{C}$, $12^\circ\text{C}$ ও $11^\circ\text{C}$ তাপমাত্রায় সম্পৃক্ত বাষ্পচাপ যথাক্রমে $22.38 \times 10^{-3}\text{ m}$, $10.52 \times 10^{-3}\text{ m}$ এবং $9.90 \times 10^{-3}\text{ m Hg}$ হলে ওই স্থানের আপেক্ষিক আর্দ্রতা কত?
    
    \vspace{0.8em}
    \begin{english}
    \noindent\textit{\color{darkgray}
    At a place, the dry bulb reading of a hygrometer is $24^\circ\text{C}$ and the dew point is $11.5^\circ\text{C}$. If Glaisher's factor at $24^\circ\text{C}$ is $1.72$, what is the wet bulb reading? Also, if the saturated vapour pressures at $24^\circ\text{C}$, $12^\circ\text{C}$, and $11^\circ\text{C}$ are $22.38 \times 10^{-3}\text{ m}$, $10.52 \times 10^{-3}\text{ m}$, and $9.90 \times 10^{-3}\text{ m Hg}$ respectively, what is the relative humidity of that place?
    }
    \end{english}
\end{tcolorbox}

\begin{itemize}[label={}, leftmargin=1cm, itemsep=0.5em]
    \item[\textbf{A)}] সিক্ত বাল্বের পাঠ $16.73^\circ\text{C}$ এবং আপেক্ষিক আর্দ্রতা $45.62\%$
    \item[\textbf{B)}] সিক্ত বাল্বের পাঠ $18.20^\circ\text{C}$ এবং আপেক্ষিক আর্দ্রতা $52.30\%$
    \item[\textbf{C)}] সিক্ত বাল্বের পাঠ $16.73^\circ\text{C}$ এবং আপেক্ষিক আর্দ্রতা $50.00\%$
    \item[\textbf{D)}] সিক্ত বাল্বের পাঠ $14.50^\circ\text{C}$ এবং আপেক্ষিক আর্দ্রতা $40.12\%$
\end{itemize}

\vspace{0.5em}
\begin{tcolorbox}[answerbox]
    \textbf{\color{success} উত্তর:} A) সিক্ত বাল্বের পাঠ $16.73^\circ\text{C}$ এবং আপেক্ষিক আর্দ্রতা $45.62\%$
\end{tcolorbox}
\vspace{1em}

\begin{tcolorbox}[solutionbox={সমাধান (Solution)}]
    \noindent
    \textbf{ধাপ ১: সিক্ত বাল্বের পাঠ নির্ণয়} \\
    গ্লাইশারের সমীকরণানুসারে শিশিরাঙ্ক $\theta_{dew} = \theta_1 - G (\theta_1 - \theta_2)$:
    \[ 11.5 = 24 - 1.72 (24 - \theta_2) \implies 1.72 (24 - \theta_2) = 12.5 \]
    \[ \implies 24 - \theta_2 = \frac{12.5}{1.72} = 7.2674 \implies \theta_2 = 24 - 7.2674 = \mathbf{16.73^\circ\text{C}} \]
    
    \vspace{0.4em}
    \noindent\textbf{ধাপ ২: আন্তঃপাতন (Interpolation) দ্বারা বাষ্পচাপ নির্ণয়} \\
    $11.5^\circ\text{C}$ শিশিরাঙ্কে উপস্থিত সম্পৃক্ত বাষ্পচাপ ($f$):
    \[ f = 9.90 \times 10^{-3} + (10.52 - 9.90) \times 10^{-3} \times 0.5 = 10.21 \times 10^{-3}\text{ m Hg} \]
    $24^\circ\text{C}$ এ সম্পৃক্ত বাষ্পচাপ $F = 22.38 \times 10^{-3}\text{ m Hg}$।
    
    \vspace{0.4em}
    \noindent\textbf{ধাপ ৩: আপেক্ষিক আর্দ্রতা হিসাব}
    \[ R = \frac{f}{F} \times 100\% = \frac{10.21 \times 10^{-3}}{22.38 \times 10^{-3}} \times 100\% = \mathbf{45.62\%} \]
\end{tcolorbox}

\vspace{2em}

% ------------------------------------------
% QUESTION 37 SECTION
% ------------------------------------------
\begin{tcolorbox}[questionbox={প্রশ্ন (Question) 7}]
    \noindent
    স্থির চাপে $4 \times 10^{-3}\text{ m}^3$ আয়তনের কোনো গ্যাসকে $0^\circ\text{C}$ থেকে $68.25^\circ\text{C}$ পর্যন্ত উত্তপ্ত করার ফলে এর আয়তন $1 \times 10^{-3}\text{ m}^3$ বৃদ্ধি পেল। এই তথ্যের ভিত্তিতে সেন্টিগ্রেড স্কেলে পরম শূন্য তাপমাত্রার মান কত?
    
    \vspace{0.8em}
    \begin{english}
    \noindent\textit{\color{darkgray}
    At constant pressure, a gas of volume $4 \times 10^{-3}\text{ m}^3$ is heated from $0^\circ\text{C}$ to $68.25^\circ\text{C}$, resulting in a volume increase of $1 \times 10^{-3}\text{ m}^3$. Based on this data, what is the value of absolute zero temperature on the Centigrade scale?
    }
    \end{english}
\end{tcolorbox}

\begin{itemize}[label={}, leftmargin=1cm, itemsep=0.5em]
    \item[\textbf{A)}] পরম শূন্য তাপমাত্রা $-273^\circ\text{C}$
    \item[\textbf{B)}] পরম শূন্য তাপমাত্রা $-268^\circ\text{C}$
    \item[\textbf{C)}] পরম শূন্য তাপমাত্রা $-280^\circ\text{C}$
    \item[\textbf{D)}] পরম শূন্য তাপমাত্রা $-270^\circ\text{C}$
\end{itemize}

\vspace{0.5em}
\begin{tcolorbox}[answerbox]
    \textbf{\color{success} উত্তর:} A) পরম শূন্য তাপমাত্রা $-273^\circ\text{C}$
\end{tcolorbox}
\vspace{1em}

\begin{tcolorbox}[solutionbox={সমাধান (Solution)}]
    \noindent
    স্থির চাপে চার্লসের সূত্রানুসারে আয়তন সম্প্রসারণ $\Delta V = V_0 \gamma_p \theta$। আয়তন প্রসারক গুণাঙ্ক:
    \[ \gamma_p = \frac{\Delta V}{V_0 \theta} = \frac{1 \times 10^{-3}}{(4 \times 10^{-3}) \times 68.25} = \frac{1}{273}\text{ }^\circ\text{C}^{-1} \]
    সেন্টিগ্রেড স্কেলে পরম শূন্য তাপমাত্রা:
    \[ T_0 = -\frac{1}{\gamma_p} = \mathbf{-273^\circ\text{C}} \]
\end{tcolorbox}

\vspace{2em}

% ------------------------------------------
% QUESTION 38 SECTION
% ------------------------------------------
\begin{tcolorbox}[questionbox={প্রশ্ন (Question) 8}]
    \noindent
    একটি ঘরের পরিমাপ $12\text{ m} \times 9\text{ m} \times 8\text{ m}$। সকালে ঘরটির তাপমাত্রা $20^\circ\text{C}$ ছিল। দুপুরে ঘরের অভ্যন্তরের তাপমাত্রা বেড়ে $32^\circ\text{C}$ হলেও বায়ুমণ্ডলীয় চাপ অপরিবর্তিত থাকে। তাপমাত্রা বৃদ্ধির ফলে প্রাথমিক বায়ুর শতকরা কত অংশ ঘর থেকে বাইরে বেরিয়ে যাবে?
    
    \vspace{0.8em}
    \begin{english}
    \noindent\textit{\color{darkgray}
    A room measures $12\text{ m} \times 9\text{ m} \times 8\text{ m}$. In the morning, the temperature inside the room was $20^\circ\text{C}$. In the afternoon, the room temperature rises to $32^\circ\text{C}$ while the atmospheric pressure remains constant. What percentage of the initial air in the room will escape to the outside due to this temperature increase?
    }
    \end{english}
\end{tcolorbox}

\begin{itemize}[label={}, leftmargin=1cm, itemsep=0.5em]
    \item[\textbf{A)}] বায়ুর $4.10\%$ বাইরে বেরিয়ে যাবে
    \item[\textbf{B)}] বায়ুর $2.50\%$ বাইরে বেরিয়ে যাবে
    \item[\textbf{C)}] বায়ুর $6.20\%$ বাইরে বেরিয়ে যাবে
    \item[\textbf{D)}] বায়ুর $8.50\%$ বাইরে বেরিয়ে যাবে
\end{itemize}

\vspace{0.5em}
\begin{tcolorbox}[answerbox]
    \textbf{\color{success} উত্তর:} A) বায়ুর $4.10\%$ বাইরে বেরিয়ে যাবে
\end{tcolorbox}
\vspace{1em}

\begin{tcolorbox}[solutionbox={সমাধান (Solution)}]
    \noindent
    ঘরের মোট আয়তন $V_1 = 12 \times 9 \times 8 = 864\text{ m}^3$। প্রাথমিক তাপমাত্রা $T_1 = 273 + 20 = 293\text{ K}$ এবং চূড়ান্ত তাপমাত্রা $T_2 = 273 + 32 = 305\text{ K}$। স্থির চাপে চার্লসের সূত্রানুসারে প্রসারিত আয়তন:
    \[ V_2 = V_1 \times \frac{T_2}{T_1} = 864 \times \frac{305}{293} = 899.385\text{ m}^3 \]
    প্রসারিত বা বেরিয়ে যাওয়া বায়ুর আয়তন:
    \[ \Delta V = V_2 - V_1 = 899.385 - 864 = 35.385\text{ m}^3 \]
    প্রাথমিক আয়তনের সাপেক্ষে শতকরা নিঃসারিত বায়ু:
    \[ \text{Percentage Escaped} = \frac{\Delta V}{V_1} \times 100\% = \frac{305 - 293}{293} \times 100\% = \mathbf{4.10\%} \]
\end{tcolorbox}

\vspace{2em}

% ------------------------------------------
% QUESTION 39 SECTION
% ------------------------------------------
\begin{tcolorbox}[questionbox={প্রশ্ন (Question) 9}]
    \noindent
    একজন ব্যক্তি প্রশ্বাসের সময় $37^\circ\text{C}$ দৈহিক তাপমাত্রায় ও $1\text{ atm}$ চাপে ফুসফুসে $5.5\text{ L}$ বাতাস গ্রহণ করতে পারে। বাতাসে আয়তন হিসেবে $21\%$ অক্সিজেন থাকলে, ওই ব্যক্তির ফুসফুসে কতটি অক্সিজেন অণু জমা হবে? ($R = 0.0821\text{ L atm mol}^{-1}\text{ K}^{-1}$, $N_A = 6.023 \times 10^{23}\text{ mol}^{-1}$)
    
    \vspace{0.8em}
    \begin{english}
    \noindent\textit{\color{darkgray}
    When a person inhales, their lungs can hold $5.5\text{ L}$ of air at a body temperature of $37^\circ\text{C}$ and $1\text{ atm}$ pressure. If air contains $21\%$ oxygen by volume, how many oxygen molecules are contained in the lungs?
    }
    \end{english}
\end{tcolorbox}

\begin{itemize}[label={}, leftmargin=1cm, itemsep=0.5em]
    \item[\textbf{A)}] অক্সিজেন অণুর সংখ্যা $2.733 \times 10^{22}$
    \item[\textbf{B)}] অক্সিজেন অণুর সংখ্যা $1.301 \times 10^{23}$
    \item[\textbf{C)}] অক্সিজেন অণুর সংখ্যা $5.420 \times 10^{21}$
    \item[\textbf{D)}] অক্সিজেন অণুর সংখ্যা $6.023 \times 10^{23}$
\end{itemize}

\vspace{0.5em}
\begin{tcolorbox}[answerbox]
    \textbf{\color{success} উত্তর:} A) অক্সিজেন অণুর সংখ্যা $2.733 \times 10^{22}$
\end{tcolorbox}
\vspace{1em}

\begin{tcolorbox}[solutionbox={সমাধান (Solution)}]
    \noindent
    ফুসফুসে গৃহীত মোট বাতাসের মোল সংখ্যা ($n_{total}$):
    \[ n_{total} = \frac{P V}{R T} = \frac{1 \times 5.5}{0.0821 \times (273 + 37)} = \frac{5.5}{0.0821 \times 310} = 0.2161\text{ mol} \]
    বাতাসে $21\%$ অক্সিজেন থাকায় অক্সিজেনের মোল সংখ্যা ($n_{\text{O}_2}$):
    \[ n_{\text{O}_2} = 0.21 \times n_{total} = 0.21 \times 0.2161 = 0.04538\text{ mol} \]
    অক্সিজেন অণুর সংখ্যা ($N_{\text{O}_2}$):
    \[ N_{\text{O}_2} = n_{\text{O}_2} \times N_A = 0.04538 \times (6.023 \times 10^{23}) = \mathbf{2.733 \times 10^{22}} \]
\end{tcolorbox}
\vspace{2em}

% ------------------------------------------
% QUESTION 57 SECTION
% ------------------------------------------
\begin{tcolorbox}[questionbox={প্রশ্ন (Question) 10}]
    \noindent
    একটি উল্লম্ব পাত্রে $1\text{ mol}$ এক-পারমাণবিক আদর্শ গ্যাস একটি $20\text{ kg}$ ভর এবং $0.02\text{ m}^2$ প্রস্থচ্ছেদের ক্ষেত্রফল বিশিষ্ট ঘর্ষণহীন পিস্টন দ্বারা আবদ্ধ রয়েছে। বায়ুমণ্ডলের চাপ $1.0 \times 10^5\text{ Pa}$। গ্যাসটির প্রসারণের সময় চাপ আয়তনের সাথে $P(V) = P_0 - \alpha V$ সম্পর্ক মেনে চলে, যেখানে $P_0 = 4.0 \times 10^5\text{ Pa}$ এবং $\alpha = 5.0 \times 10^7\text{ Pa m}^{-3}$। আয়তন $V_1 = 0.001\text{ m}^3$ হতে $V_2 = 0.003\text{ m}^3$ এ প্রসারিত হলে যোগজীকরণ প্রক্রিয়ায় পিস্টনের ওপর লব্ধি যান্ত্রিক কাজ এবং পিস্টনের বেগ সর্বোচ্চ হওয়ার মুহূর্তে গ্যাসের আয়তন কত হবে? ($g = 9.8\text{ m s}^{-2}$)
    
    \vspace{0.8em}
    \begin{english}
    \noindent\textit{\color{darkgray}
    A vertical cylinder contains $1\text{ mol}$ of a monatomic ideal gas enclosed by a frictionless piston of mass $20\text{ kg}$ and cross-sectional area $0.02\text{ m}^2$. The atmospheric pressure is $1.0 \times 10^5\text{ Pa}$. During expansion, the gas pressure varies with volume as $P(V) = P_0 - \alpha V$, where $P_0 = 4.0 \times 10^5\text{ Pa}$ and $\alpha = 5.0 \times 10^7\text{ Pa m}^{-3}$. If the volume expands from $V_1 = 0.001\text{ m}^3$ to $V_2 = 0.003\text{ m}^3$, calculate using integration the net mechanical work done on the piston and determine using differentiation the volume at which the piston reaches maximum speed. ($g = 9.8\text{ m s}^{-2}$)
    }
    \end{english}
\end{tcolorbox}

\begin{itemize}[label={}, leftmargin=1cm, itemsep=0.5em]
    \item[\textbf{A)}] লব্ধি কাজ $380.4\text{ J}$ এবং আয়তন $0.005804\text{ m}^3$
    \item[\textbf{B)}] লব্ধি কাজ $600.0\text{ J}$ এবং আয়তন $0.003000\text{ m}^3$
    \item[\textbf{C)}] লব্ধি কাজ $200.0\text{ J}$ এবং আয়তন $0.001000\text{ m}^3$
    \item[\textbf{D)}] লব্ধি কাজ $450.2\text{ J}$ এবং আয়তন $0.008200\text{ m}^3$
\end{itemize}

\vspace{0.5em}
\begin{tcolorbox}[answerbox]
    \textbf{\color{success} উত্তর:} A) লব্ধি কাজ $380.4\text{ J}$ এবং আয়তন $0.005804\text{ m}^3$
\end{tcolorbox}
\vspace{1em}

\begin{tcolorbox}[solutionbox={সমাধান (Solution)}]
    \noindent
    \textbf{ধাপ ১: যোগজীকরণ প্রক্রিয়ায় গ্যাস দ্বারা পিস্টনের ওপর সম্পাদিত কাজ}
    \begin{align*}
        W_{gas} &= \int_{V_1}^{V_2} P(V) \, dV = \int_{0.001}^{0.003} (4 \times 10^5 - 5 \times 10^7 V) \, dV \\
        &= \left[ 4 \times 10^5 V - \frac{5 \times 10^7}{2} V^2 \right]_{0.001}^{0.003} \\
        &= 4 \times 10^5 \times (0.003 - 0.001) - 2.5 \times 10^7 \times (0.003^2 - 0.001^2) \\
        &= 800 - 2.5 \times 10^7 \times (8 \times 10^{-6}) = 800 - 200 = 600\text{ J}
    \end{align*}
    
    \vspace{0.4em}
    \noindent\textbf{ধাপ ২: বায়ুমণ্ডলীয় চাপ ও মহাকর্ষের বিরুদ্ধে কাজ হিসাব}
    \begin{itemize}[leftmargin=1.5em, itemsep=0.2em]
        \item বায়ুমণ্ডলীয় চাপের বিরুদ্ধে কাজ, $W_{atm} = P_{atm} \Delta V = 1.0 \times 10^5 \times (0.003 - 0.001) = 200\text{ J}$
        \item পিস্টনের সরণ $\Delta h = \frac{\Delta V}{A} = \frac{0.002}{0.02} = 0.1\text{ m}$
        \item মহাকর্ষের বিরুদ্ধে কাজ,  $W_{gravity} = m g \Delta h = 20 \times 9.8 \times 0.1 = 19.6\text{ J}$
    \end{itemize}
    পিস্টনের ওপর লব্ধি কাজ:
    \[ W_{net} = W_{gas} - W_{atm} - W_{gravity} = 600 - 200 - 19.6 = \mathbf{380.4\text{ J}} \]
    
    \vspace{0.4em}
    \noindent\textbf{ধাপ ৩: পিস্টনের বেগ সর্বোচ্চ হওয়ার শর্ত প্রয়োগ}
    পিস্টনের বেগ সর্বোচ্চ হবে যখন লব্ধি বল $F_{net} = 0$ হবে:
    \[ F_{net}(V) = [P(V) - P_{atm}] A - m g = 0 \]
    \[ [(4 \times 10^5 - 5 \times 10^7 V) - 1.0 \times 10^5] \times 0.02 - 20 \times 9.8 = 0 \]
    \[ (3 \times 10^5 - 5 \times 10^7 V) \times 0.02 = 196 \implies 6000 - 10^6 V = 196 \]
    \[ 10^6 V = 5804 \implies V = \mathbf{0.005804\text{ m}^3} \]
\end{tcolorbox}

\vspace{2em}

% ------------------------------------------
% QUESTION 75 SECTION
% ------------------------------------------
\begin{tcolorbox}[questionbox={প্রশ্ন (Question) 11}]
    \noindent
    একটি ঘরের তাপমাত্রা $30^\circ\text{C}$ এবং আপেক্ষিক আর্দ্রতা $70\%$। $30^\circ\text{C}$ তাপমাত্রায় সম্পৃক্ত বাষ্পচাপ $31.8\text{ mm Hg}$। যদি ঘরের তাপমাত্রা কমিয়ে $20^\circ\text{C}$ করা হয়, যেখানে সম্পৃক্ত বাষ্পচাপ $17.5\text{ mm Hg}$, তবে প্রতি ঘনমিটারে ঘনীভূত পানির ভর কত? ($R = 8.314\text{ J mol}^{-1}\text{ K}^{-1}$, পানির আণবিক ভর $M = 18\text{ g mol}^{-1}$)
    
    \vspace{0.8em}
    \begin{english}
    \noindent\textit{\color{darkgray}
    The room temperature is $30^\circ\text{C}$ and its relative humidity is $70\%$. The saturated vapour pressure at $30^\circ\text{C}$ is $31.8\text{ mm Hg}$. If the room temperature is lowered to $20^\circ\text{C}$, where the saturated vapour pressure is $17.5\text{ mm Hg}$, what mass of water will condense per cubic meter? ($R = 8.314\text{ J mol}^{-1}\text{ K}^{-1}$, $M = 18\text{ g mol}^{-1}$)
    }
    \end{english}
\end{tcolorbox}

\begin{itemize}[label={}, leftmargin=1cm, itemsep=0.5em]
    \item[\textbf{A)}] $3.97\text{ g m}^{-3}$
    \item[\textbf{B)}] $14.30\text{ g m}^{-3}$
    \item[\textbf{C)}] $8.50\text{ g m}^{-3}$
    \item[\textbf{D)}] $1.25\text{ g m}^{-3}$
\end{itemize}

\vspace{0.5em}
\begin{tcolorbox}[answerbox]
    \textbf{\color{success} উত্তর:} A) $3.97\text{ g m}^{-3}$
\end{tcolorbox}
\vspace{1em}

\begin{tcolorbox}[solutionbox={সমাধান (Solution)}]
    \noindent
    \textbf{ধাপ ১: প্রাথমিক বাষ্পচাপ ও ঘনত্ব নির্ণয়} \\
    $30^\circ\text{C}$ ($303\text{ K}$) তাপমাত্রায় বাতাসে উপস্থিত বাষ্পচাপ:
    \[ f_1 = 0.70 \times 31.8\text{ mm Hg} = 22.26\text{ mm Hg} = 2967.75\text{ Pa} \]
    বাতাসে জলীয় বাষ্পের প্রাথমিক ঘনত্ব:
    \[ \rho_1 = \frac{f_1 M}{R T_1} = \frac{2967.75 \times 0.018}{8.314 \times 303} = 0.021205\text{ kg m}^{-3} = 21.205\text{ g m}^{-3} \]
    
    \vspace{0.4em}
    \noindent\textbf{ধাপ ২: অন্তিম ঘনত্ব ও ঘনীভূত পানির ভর হিসাব} \\
    $20^\circ\text{C}$ ($293\text{ K}$) তাপমাত্রায় সম্পৃক্ত বাষ্পচাপ $F_2 = 17.5\text{ mm Hg} = 2333.14\text{ Pa}$। বায়ুতে অবশিষ্ট বাষ্পের সর্বোচ্চ ঘনত্ব:
    \[ \rho_2 = \frac{F_2 M}{R T_2} = \frac{2333.14 \times 0.018}{8.314 \times 293} = 0.017240\text{ kg m}^{-3} = 17.240\text{ g m}^{-3} \]
    প্রতি ঘনমিটারে ঘনীভূত পানির ভর:
    \[ \Delta \rho = \rho_1 - \rho_2 = 21.205 - 17.240 = 3.965\text{ g m}^{-3} \approx \mathbf{3.97\text{ g m}^{-3}} \]
\end{tcolorbox}

\vspace{2em}

% ------------------------------------------
% QUESTION 76 SECTION
% ------------------------------------------
\begin{tcolorbox}[questionbox={প্রশ্ন (Question) 12}]
    \noindent
    $1\text{ mol}$ এক-পারমাণবিক আদর্শ গ্যাস ($M = 4\text{ g mol}^{-1}$) একটি প্রক্রিয়ার মধ্য দিয়ে প্রসারিত হয় যেখানে চাপ আয়তনের সাথে $P(V) = P_0 (1 - \beta V)$ নিয়ম মেনে চলে। $P_0 = 1.0 \times 10^5\text{ Pa}$ এবং $\beta = 100\text{ m}^{-3}$ হলে, এই প্রসারণ প্রক্রিয়ায় অণুগুলোর অর্জিত সর্বোচ্চ মূল গড় বর্গবেগ ($c_{\text{rms,max}}$) কত হবে? ($R = 8.314\text{ J mol}^{-1}\text{ K}^{-1}$)
    
    \vspace{0.8em}
    \begin{english}
    \noindent\textit{\color{darkgray}
    $1\text{ mol}$ of a monatomic ideal gas ($M = 4\text{ g mol}^{-1}$) expands in a process where pressure varies with volume as $P(V) = P_0 (1 - \beta V)$. If $P_0 = 1.0 \times 10^5\text{ Pa}$ and $\beta = 100\text{ m}^{-3}$, what will be the maximum root mean square speed ($c_{\text{rms,max}}$) reached by the gas molecules during this expansion? ($R = 8.314\text{ J mol}^{-1}\text{ K}^{-1}$)
    }
    \end{english}
\end{tcolorbox}

\begin{itemize}[label={}, leftmargin=1cm, itemsep=0.5em]
    \item[\textbf{A)}] $433.01\text{ m s}^{-1}$
    \item[\textbf{B)}] $250.00\text{ m s}^{-1}$
    \item[\textbf{C)}] $612.37\text{ m s}^{-1}$
    \item[\textbf{D)}] $300.00\text{ m s}^{-1}$
\end{itemize}

\vspace{0.5em}
\begin{tcolorbox}[answerbox]
    \textbf{\color{success} উত্তর:} A) $433.01\text{ m s}^{-1}$
\end{tcolorbox}
\vspace{1em}

\begin{tcolorbox}[solutionbox={সমাধান (Solution)}]
    \noindent
    \textbf{ধাপ ১: তাপমাত্রা ও সর্বোচ্চ তাপমাত্রার শর্ত} \\
    আদর্শ গ্যাস সমীকরণ $PV = nRT$ হতে তাপমাত্রা:
    \[ T(V) = \frac{P(V) V}{n R} = \frac{P_0 (V - \beta V^2)}{n R} \]
    তাপমাত্রা সর্বোচ্চ হওয়ার জন্য আয়তনের সাপেক্ষে ব্যবকলন (differentiation) করে শূন্যের সমান ধরি:
    \[ \frac{dT}{dV} = \frac{P_0}{n R} (1 - 2 \beta V) = 0 \implies V = \frac{1}{2 \beta} \]
    মান বসিয়ে সর্বোচ্চ তাপমাত্রা $T_{max}$:
    \[ T_{max} = \frac{P_0}{4 n R \beta} = \frac{1.0 \times 10^5}{4 \times 1 \times 8.314 \times 100} = 30.07\text{ K} \]
    
    \vspace{0.4em}
    \noindent\textbf{ধাপ ২: সর্বোচ্চ মূল গড় বর্গবেগ হিসাব}
    \[ c_{rms, max} = \sqrt{\frac{3 R T_{max}}{M}} = \sqrt{\frac{3 \times 8.314 \times 30.07}{4 \times 10^{-3}}} = \sqrt{187500} = \mathbf{433.01\text{ m s}^{-1}} \]
\end{tcolorbox}
% ------------------------------------------
% QUESTION 13 SECTION
% ------------------------------------------
\begin{tcolorbox}[questionbox={প্রশ্ন (Question) 13}]
    \noindent
    $0.2\text{ m}$ ব্যাসার্ধের একটি সিলিন্ডার আকৃতির পাত্রে $300\text{ K}$ তাপমাত্রায় নাইট্রোজেন গ্যাস ($M = 28\text{ g mol}^{-1}$) আবদ্ধ আছে। সিলিন্ডারটি তার কেন্দ্রীয় অক্ষের চারপাশে $1000\text{ rad s}^{-1}$ কৌণিক বেগে ঘুরতে থাকলে কেন্দ্রিয় অক্ষের চাপ $P_0$ এবং সিলিন্ডারের বাইরের দেয়ালের চাপ $P(R_0)$ এর অনুপাত $\frac{P(R_0)}{P_0}$ কত হবে? এবং কেন্দ্রিয় অক্ষের চাপ $P_0 = 1.0 \times 10^5\text{ Pa}$ হলে দেয়ালের গায়ে চাপের পরিবর্তনের নতি বা নতিমাত্রা $\frac{dP}{dr}$ কত হবে? ($R = 8.314\text{ J mol}^{-1}\text{ K}^{-1}$)
    
    \vspace{0.8em}
    \begin{english}
    \noindent\textit{\color{darkgray}
    A cylindrical container of radius $0.2\text{ m}$ contains Nitrogen gas ($M = 28\text{ g mol}^{-1}$) at $300\text{ K}$. If the cylinder rotates about its central axis at an angular velocity of $1000\text{ rad s}^{-1}$, what will be the ratio of pressure at the outer wall to the central axis $\frac{P(R_0)}{P_0}$? And if the central axis pressure is $P_0 = 1.0 \times 10^5\text{ Pa}$, what will be the pressure gradient $\frac{dP}{dr}$ at the outer wall? ($R = 8.314\text{ J mol}^{-1}\text{ K}^{-1}$)
    }
    \end{english}
\end{tcolorbox}

% ------------------------------------------
% TIKZ DIAGRAM FOR ROTATING CYLINDER
% ------------------------------------------
\vspace{1em}
\begin{center}
\begin{tikzpicture}[scale=1.1, >=Stealth]
    % Cylinder Top View / Cross-section representation
    \draw[thick, darkgray, fill=secondary!30] (0, 0) circle (2.2cm);
    
    % Central Axis (Point)
    \fill[primary] (0, 0) circle (2.5pt);
    \node[above, font=\small, text=primary] at (0, 0.1) {অক্ষ ($P_0$)};
    
    % Radius line to outer wall
    \draw[->, thick, accent] (0, 0) -- (1.55, 1.55) node[midway, above left, font=\small] {$R_0 = 0.2\text{ m}$};
    
    % Rotation Arrow
    \draw[->, thick, primary, domain=0:270, samples=50] plot ({1.8*cos(\x)}, {1.8*sin(\x)});
    \node[right, font=\small, text=primary] at (1.8, 0) {$\omega = 1000\text{ rad/s}$};
    
    % Outer Wall Pressure Label
    \node[font=\small, text=darkgray] at (0, -2.5) {বাইরের দেওয়াল: $P(R_0)$};

\end{tikzpicture}
\end{center}
\vspace{1em}

\begin{itemize}[label={}, leftmargin=1cm, itemsep=0.5em]
    \item[\textbf{A)}] চাপের অনুপাত $1.252$ এবং চাপের নতিমাত্রা $2.810 \times 10^5\text{ Pa m}^{-1}$
    \item[\textbf{B)}] চাপের অনুপাত $2.500$ এবং চাপের নতিমাত্রা $5.400 \times 10^5\text{ Pa m}^{-1}$
    \item[\textbf{C)}] চাপের অনুপাত $1.050$ এবং চাপের নতিমাত্রা $1.200 \times 10^5\text{ Pa m}^{-1}$
    \item[\textbf{D)}] চাপের অনুপাত $1.800$ এবং চাপের নতিমাত্রা $3.500 \times 10^5\text{ Pa m}^{-1}$
\end{itemize}

\vspace{0.5em}
\begin{tcolorbox}[answerbox]
    \textbf{\color{success} উত্তর:} A) চাপের অনুপাত $1.252$ এবং চাপের নতিমাত্রা $2.810 \times 10^5\text{ Pa m}^{-1}$
\end{tcolorbox}
\vspace{1em}

\begin{tcolorbox}[solutionbox={সমাধান (Solution)}]
    \noindent
    \textbf{ধাপ ১: ঘূর্ণায়মান পাত্রের চাপের সমীকরণ} \\
    ঘূর্ণায়মান গ্যাসের ক্ষেত্রে ব্যাসার্ধ $r$ বরাবর চাপের পরিবর্তন:
    \[ \frac{dP/dr}{P} = \frac{M \omega^2 r}{R T} \implies P(r) = P_0 e^{\frac{M \omega^2 r^2}{2 R T}} \]
    
    \vspace{0.4em}
    \noindent\textbf{ধাপ ২: দেয়াল ও অক্ষের চাপের অনুপাত নির্ণয়}
    \begin{itemize}[leftmargin=1.5em, itemsep=0.2em]
        \item মোলার ভর, $M = 28 \times 10^{-3}\text{ kg mol}^{-1}$
        \item তাপমাত্রা, $T = 300\text{ K}$
        \item কৌণিক বেগ, $\omega = 1000\text{ rad s}^{-1}$
        \item ব্যাসার্ধ, $R_0 = 0.2\text{ m}$
    \end{itemize}
    ঘাতের মান:
    \[ \frac{M \omega^2 R_0^2}{2 R T} = \frac{(28 \times 10^{-3}) \times (1000)^2 \times (0.2)^2}{2 \times 8.314 \times 300} = \frac{1120}{4988.4} = 0.2245 \]
    চাপের অনুপাত:
    \[ \frac{P(R_0)}{P_0} = e^{0.2245} = \mathbf{1.252} \]
    
    \vspace{0.4em}
    \noindent\textbf{ধাপ ৩: দেয়ালের গায়ে চাপের নতিমাত্রা ($\frac{dP}{dr}$) নির্ণয়} \\
    দেওয়া আছে $P_0 = 1.0 \times 10^5\text{ Pa}$, ফলে বাইরের দেয়ালে চাপ $P(R_0) = 1.252 \times 1.0 \times 10^5 = 1.252 \times 10^5\text{ Pa}$। চাপের নতিমাত্রা:
    \begin{align*}
        \frac{dP}{dr}\Big|_{r=R_0} &= P(R_0) \times \frac{M \omega^2 R_0}{R T} \\
        &= 1.252 \times 10^5 \times \frac{(28 \times 10^{-3}) \times (1000)^2 \times 0.2}{8.314 \times 300} \\
        &= 1.252 \times 10^5 \times \frac{5600}{2494.2} \\
        &= 1.252 \times 10^5 \times 2.2452 = \mathbf{2.810 \times 10^5\text{ Pa m}^{-1}}
    \end{align*}
\end{tcolorbox}

\vspace{2em}

% ------------------------------------------
% QUESTION 14 SECTION
% ------------------------------------------
\begin{tcolorbox}[questionbox={প্রশ্ন (Question) 14}]
    \noindent
    নির্দিষ্ট তাপমাত্রায় $T$ একটি আবদ্ধ পাত্রে আদর্শ গ্যাসের অণুগুলোর $c_{\text{rms}}$, $c_{\text{avg}}$ এবং $c_{\text{mp}}$ বেগের অনুপাত কত? যদি সমতাপীয় প্রক্রিয়ায় গ্যাসের আয়তন দ্বিগুণ করা হয়, তবে অণুগুলোর মোট গতিশক্তি এবং $c_{\text{rms}} : c_{\text{avg}} : c_{\text{mp}}$ অনুপাতের কী পরিবর্তন ঘটবে?
    
    \vspace{0.8em}
    \begin{english}
    \noindent\textit{\color{darkgray}
    For an ideal gas at a given temperature $T$, what is the ratio of $c_{\text{rms}}$, $c_{\text{avg}}$, and $c_{\text{mp}}$? If the gas undergoes isothermal expansion such that its volume doubles, how will the total kinetic energy and the ratio $c_{\text{rms}} : c_{\text{avg}} : c_{\text{mp}}$ change?
    }
    \end{english}
\end{tcolorbox}

\begin{itemize}[label={}, leftmargin=1cm, itemsep=0.5em]
    \item[\textbf{A)}] $\sqrt{3} : \sqrt{\frac{8}{\pi}} : \sqrt{2}$; Unchanged
    \item[\textbf{B)}] $1 : 1 : 1$; Doubles
    \item[\textbf{C)}] $\sqrt{2} : \sqrt{\frac{8}{\pi}} : \sqrt{3}$; Halved
    \item[\textbf{D)}] $1.5 : 1.2 : 1.0$; Quadrupled
\end{itemize}

\vspace{0.5em}
\begin{tcolorbox}[answerbox]
    \textbf{\color{success} উত্তর:} A) $\sqrt{3} : \sqrt{\frac{8}{\pi}} : \sqrt{2}$; Unchanged
\end{tcolorbox}
\vspace{1em}

\begin{tcolorbox}[solutionbox={সমাধান (Solution)}]
    \noindent
    আদর্শ গ্যাসের আণবিক বেগসমূহের অনুপাত:
    \[ c_{\text{rms}} : c_{\text{avg}} : c_{\text{mp}} = \sqrt{\frac{3RT}{M}} : \sqrt{\frac{8RT}{\pi M}} : \sqrt{\frac{2RT}{M}} = \mathbf{\sqrt{3} : \sqrt{\frac{8}{\pi}} : \sqrt{2}} \approx 1.225 : 1.128 : 1 \]
    যেহেতু প্রক্রিয়াটি সমতাপীয় (isothermal), তাই পরম তাপমাত্রা $T$ স্থির থাকে। আদর্শ গ্যাসের মোট গতিশক্তি $E_k = \frac{f}{2} n R T$ শুধুমাত্র তাপমাত্রার ওপর নির্ভর করে, তাই আয়তন দ্বিগুণ হলেও গতিশক্তি এবং বেগের অনুপাত উভয়ই **অপরিবর্তিত** থাকবে।
\end{tcolorbox}

\vspace{2em}

% ------------------------------------------
% QUESTION 15 SECTION
% ------------------------------------------
\begin{tcolorbox}[questionbox={প্রশ্ন (Question) 15}]
    \noindent
    $1.0 \times 10^5\text{ Pa}$ চাপে এবং নির্দিষ্ট তাপমাত্রায় একটি আদর্শ গ্যাসের ঘনত্ব $1.2\text{ kg m}^{-3}$। গ্যাসটির অণুগুলোর মূল গড় বর্গবেগ ($c_{\text{rms}}$), গড় বেগ ($c_{\text{avg}}$) এবং সর্বাধিক সম্ভাব্য বেগ ($c_{\text{mp}}$) এর মান কত?
    
    \vspace{0.8em}
    \begin{english}
    \noindent\textit{\color{darkgray}
    The density of an ideal gas is $1.2\text{ kg m}^{-3}$ at a pressure of $1.0 \times 10^5\text{ Pa}$ and a certain temperature. What are the values of the root mean square speed ($c_{\text{rms}}$), average speed ($c_{\text{avg}}$), and most probable speed ($c_{\text{mp}}$) of the gas molecules?
    }
    \end{english}
\end{tcolorbox}

\begin{itemize}[label={}, leftmargin=1cm, itemsep=0.5em]
    \item[\textbf{A)}] $c_{\text{rms}} = 500.00\text{ m s}^{-1}, c_{\text{avg}} = 460.66\text{ m s}^{-1}, c_{\text{mp}} = 408.25\text{ m s}^{-1}$
    \item[\textbf{B)}] $c_{\text{rms}} = 408.25\text{ m s}^{-1}, c_{\text{avg}} = 500.00\text{ m s}^{-1}, c_{\text{mp}} = 460.66\text{ m s}^{-1}$
    \item[\textbf{C)}] $c_{\text{rms}} = 600.00\text{ m s}^{-1}, c_{\text{avg}} = 520.00\text{ m s}^{-1}, c_{\text{mp}} = 450.00\text{ m s}^{-1}$
    \item[\textbf{D)}] $c_{\text{rms}} = 350.50\text{ m s}^{-1}, c_{\text{avg}} = 300.00\text{ m s}^{-1}, c_{\text{mp}} = 250.00\text{ m s}^{-1}$
\end{itemize}

\vspace{0.5em}
\begin{tcolorbox}[answerbox]
    \textbf{\color{success} উত্তর:} A) $c_{\text{rms}} = 500.00\text{ m s}^{-1}, c_{\text{avg}} = 460.66\text{ m s}^{-1}, c_{\text{mp}} = 408.25\text{ m s}^{-1}$
\end{tcolorbox}
\vspace{1em}

\begin{tcolorbox}[solutionbox={সমাধান (Solution)}]
    \noindent
    প্রদত্ত উপাত্ত: $P = 1.0 \times 10^5\text{ Pa}$ এবং $\rho = 1.2\text{ kg m}^{-3}$।
    \begin{itemize}[leftmargin=1.5em, itemsep=0.2em]
        \item মূল গড় বর্গবেগ ($c_{\text{rms}}$):
        \[ c_{\text{rms}} = \sqrt{\frac{3 P}{\rho}} = \sqrt{\frac{3 \times 1.0 \times 10^5}{1.2}} = \sqrt{250000} = \mathbf{500.00\text{ m s}^{-1}} \]
        \item গড় বেগ ($c_{\text{avg}}$):
        \[ c_{\text{avg}} = \sqrt{\frac{8 P}{\pi \rho}} = \sqrt{\frac{8 \times 1.0 \times 10^5}{3.14159 \times 1.2}} = \sqrt{212206.6} = \mathbf{460.66\text{ m s}^{-1}} \]
        \item সর্বাধিক সম্ভাব্য বেগ ($c_{\text{mp}}$):
        \[ c_{\text{mp}} = \sqrt{\frac{2 P}{\rho}} = \sqrt{\frac{2 \times 1.0 \times 10^5}{1.2}} = \sqrt{166666.67} = \mathbf{408.25\text{ m s}^{-1}} \]
    \end{itemize}
\end{tcolorbox}

% ------------------------------------------
% QUESTION 63 SECTION
% ------------------------------------------
\begin{tcolorbox}[questionbox={প্রশ্ন (Question) 16}]
    \noindent
    আদর্শ গ্যাসের গতি তত্ত্বের মৌলিক স্বীকার্য সংক্রান্ত নিচের কোনটি সঠিক নয়?
    
    \vspace{0.8em}
    \begin{english}
    \noindent\textit{\color{darkgray}
    Which of the following statements regarding the basic postulates of the kinetic theory of ideal gases is incorrect?
    }
    \end{english}
\end{tcolorbox}

\begin{itemize}[label={}, leftmargin=1cm, itemsep=0.5em]
    \item[\textbf{A)}] Intermolecular force = 0
    \item[\textbf{B)}] Volume of molecules << Volume of container
    \item[\textbf{C)}] Collisions are perfectly inelastic
    \item[\textbf{D)}] Time of collision << Time between collisions
\end{itemize}

\vspace{0.5em}
\begin{tcolorbox}[answerbox]
    \textbf{\color{success} উত্তর:} C) Collisions are perfectly inelastic
\end{tcolorbox}
\vspace{1em}

\begin{tcolorbox}[solutionbox={সমাধান (Solution)}]
    \noindent
    আদর্শ গ্যাসের গতি তত্ত্বের মৌলিক স্বীকার্যানুসারে, গ্যাসের অণুগুলোর মধ্যে সংঘটিত সমস্ত সংঘাত সম্পূর্ণরূপে **স্থিতিস্থাপক (perfectly elastic)**। এর ফলে সংঘাতের পূর্বে ও পরে অণুগুলোর মোট রৈখিক ভরবেগ এবং মোট গতিশক্তি সংরক্ষিত থাকে। সুতরাং সংঘাত অস্থিতিস্থাপক হওয়ার দাবিটি ভুল।
\end{tcolorbox}

\vspace{2em}

% ------------------------------------------
% QUESTION 64 SECTION
% ------------------------------------------
\begin{tcolorbox}[questionbox={প্রশ্ন (Question) 17}]
    \noindent
    $27^\circ\text{C}$ তাপমাত্রায় একটি সিলিন্ডারে $20\text{ atm}$ চাপে আদর্শ গ্যাস রাখা আছে। যদি সিলিন্ডারের নিরাপত্তা ভালভটি $30\text{ atm}$ চাপে খুলে যায়, তবে বিস্ফোরণ এড়াতে সিলিন্ডারটিকে সর্বোচ্চ কত তাপমাত্রায় উত্তপ্ত করা যাবে?
    
    \vspace{0.8em}
    \begin{english}
    \noindent\textit{\color{darkgray}
    A cylinder contains an ideal gas at 20 atm pressure and 27°C temperature. If the safety valve of the cylinder opens at a pressure exceeding 30 atm, what is the maximum temperature to which the cylinder can be heated safely?
    }
    \end{english}
\end{tcolorbox}

\begin{itemize}[label={}, leftmargin=1cm, itemsep=0.5em]
    \item[\textbf{A)}] $177^\circ\text{C}$ ($450\text{ K}$)
    \item[\textbf{B)}] $200^\circ\text{C}$ ($473\text{ K}$)
    \item[\textbf{C)}] $150^\circ\text{C}$ ($423\text{ K}$)
    \item[\textbf{D)}] $300^\circ\text{C}$ ($573\text{ K}$)
\end{itemize}

\vspace{0.5em}
\begin{tcolorbox}[answerbox]
    \textbf{\color{success} উত্তর:} A) $177^\circ\text{C}$ ($450\text{ K}$)
\end{tcolorbox}
\vspace{1em}

\begin{tcolorbox}[solutionbox={সমাধান (Solution)}]
    \noindent
    সিলিন্ডারের আয়তন স্থির থাকায় গেই-লুসাকের চাপ সূত্র প্রযোজ্য:
    \[ \frac{P_1}{T_1} = \frac{P_2}{T_2} \]
    প্রদত্ত মান: $P_1 = 20\text{ atm}$, $T_1 = 273 + 27 = 300\text{ K}$ এবং $P_2 = 30\text{ atm}$। সর্বোচ্চ সহনশীল তাপমাত্রা $T_2$:
    \[ T_2 = T_1 \times \frac{P_2}{P_1} = 300 \times \frac{30}{20} = 450\text{ K} \]
    সেলসিয়াস স্কেলে:
    \[ t_2 = 450 - 273 = \mathbf{177^\circ\text{C}} \]
\end{tcolorbox}

\vspace{2em}

% ------------------------------------------
% QUESTION 65 SECTION
% ------------------------------------------
% ------------------------------------------
% QUESTION 18 SECTION
% ------------------------------------------
\begin{tcolorbox}[questionbox={প্রশ্ন (Question) 18}]
    \noindent
    $27^\circ\text{C}$ তাপমাত্রায় $4\text{ mol}$ অ-সরলরৈখিক বহুপারমাণবিক গ্যাসের (স্বাধীনতার মাত্রা $f = 6$) মোট গতিশক্তি কত? ($R = 8.314\text{ J mol}^{-1}\text{ K}^{-1}$)
    
    \vspace{0.8em}
    \begin{english}
    \noindent\textit{\color{darkgray}
    What is the total kinetic energy of $4\text{ mol}$ of a non-linear polyatomic gas (degrees of freedom $f = 6$) at $27^\circ\text{C}$? ($R = 8.314\text{ J mol}^{-1}\text{ K}^{-1}$)
    }
    \end{english}
\end{tcolorbox}

\begin{itemize}[label={}, leftmargin=1cm, itemsep=0.5em]
    \item[\textbf{A)}] $29.93\text{ kJ}$
    \item[\textbf{B)}] $14.96\text{ kJ}$
    \item[\textbf{C)}] $44.90\text{ kJ}$
    \item[\textbf{D)}] $9.98\text{ kJ}$
\end{itemize}

\vspace{0.5em}
\begin{tcolorbox}[answerbox]
    \textbf{\color{success} উত্তর:} A) $29.93\text{ kJ}$
\end{tcolorbox}
\vspace{1em}

\begin{tcolorbox}[solutionbox={সমাধান (Solution)}]
    \noindent
    সমশক্তি বণ্টন নীতিানুসারে $n$ মোল গ্যাসের মোট গতিশক্তি:
    \[ E_k = \frac{f}{2} n R T \]
    এখানে $f = 6$, $n = 4\text{ mol}$, $T = 273 + 27 = 300\text{ K}$, $R = 8.314\text{ J mol}^{-1}\text{ K}^{-1}$।
    \begin{align*}
        E_k &= \frac{6}{2} \times 4 \times 8.314 \times 300 = 3 \times 4 \times 8.314 \times 300 \\
        &= 29930.4\text{ J} = \mathbf{29.93\text{ kJ}}
    \end{align*}
\end{tcolorbox}

\vspace{2em}

% ------------------------------------------
% QUESTION 19 SECTION
% ------------------------------------------
\begin{tcolorbox}[questionbox={প্রশ্ন (Question) 19}]
    \noindent
    স্থির চাপে (constant $P$) এবং স্থির আয়তনে (constant $V$) একটি আদর্শ গ্যাসের অণুগুলোর গড় মুক্ত পথ ($\lambda$) পরম তাপমাত্রা $T$ এর সাথে কীভাবে পরিবর্তিত হয়?
    
    \vspace{0.8em}
    \begin{english}
    \noindent\textit{\color{darkgray}
    How does the mean free path ($\lambda$) of ideal gas molecules vary with absolute temperature $T$ at constant pressure $P$, and at constant volume $V$?
    }
    \end{english}
\end{tcolorbox}

\begin{itemize}[label={}, leftmargin=1cm, itemsep=0.5em]
    \item[\textbf{A)}] $P = \text{const} \implies \lambda \propto T; \quad V = \text{const} \implies \lambda = \text{const}$
    \item[\textbf{B)}] $P = \text{const} \implies \lambda \propto \frac{1}{T}; \quad V = \text{const} \implies \lambda \propto T$
    \item[\textbf{C)}] $\lambda \propto \sqrt{T}$
    \item[\textbf{D)}] $\lambda \propto T^2$
\end{itemize}

\vspace{0.5em}
\begin{tcolorbox}[answerbox]
    \textbf{\color{success} উত্তর:} A) $P = \text{const} \implies \lambda \propto T; \quad V = \text{const} \implies \lambda = \text{const}$
\end{tcolorbox}
\vspace{1em}

\begin{tcolorbox}[solutionbox={সমাধান (Solution)}]
    \noindent
    ম্যাক্সওয়েলের গড় মুক্ত পথের সমীকরণ:
    \[ \lambda = \frac{1}{\sqrt{2} \pi \sigma^2 n_V} = \frac{k T}{\sqrt{2} \pi \sigma^2 P} \]
    ১. স্থির চাপে ($P = \text{const}$): $\lambda \propto T$ (তাপমাত্রা বাড়লে গড় মুক্ত পথ সমানুপাতিক হারে বাড়ে)। \\
    ২. স্থির আয়তনে ($V = \text{const}$): প্রতি একক আয়তনে অণুর সংখ্যা $n_V = \frac{N}{V}$ স্থির থাকে, তাই $\lambda$ তাপমাত্রার ওপর নির্ভর করে না (ধ্রুবক থাকে)।
\end{tcolorbox}

\vspace{2em}

% ------------------------------------------
% QUESTION 20 SECTION
% ------------------------------------------
\begin{tcolorbox}[questionbox={প্রশ্ন (Question) 20}]
    \noindent
    একটি পাত্রে সমান ভরের নাইট্রোজেন ($M_1 = 28\text{ g mol}^{-1}$) এবং অক্সিজেন ($M_2 = 32\text{ g mol}^{-1}$) গ্যাসের মিশ্রণ রয়েছে। মিশ্রণের মোট চাপ $300\text{ kPa}$ হলে নাইট্রোজেন গ্যাসের আংশিক চাপ কত?
    
    \vspace{0.8em}
    \begin{english}
    \noindent\textit{\color{darkgray}
    A vessel contains equal masses of Nitrogen ($M_1 = 28\text{ g mol}^{-1}$) and Oxygen ($M_2 = 32\text{ g mol}^{-1}$). If the total pressure of the mixture is $300\text{ kPa}$, what is the partial pressure of Nitrogen gas?
    }
    \end{english}
\end{tcolorbox}

\begin{itemize}[label={}, leftmargin=1cm, itemsep=0.5em]
    \item[\textbf{A)}] $160\text{ kPa}$
    \item[\textbf{B)}] $140\text{ kPa}$
    \item[\textbf{C)}] $150\text{ kPa}$
    \item[\textbf{D)}] $180\text{ kPa}$
\end{itemize}

\vspace{0.5em}
\begin{tcolorbox}[answerbox]
    \textbf{\color{success} উত্তর:} A) $160\text{ kPa}$
\end{tcolorbox}
\vspace{1em}

\begin{tcolorbox}[solutionbox={সমাধান (Solution)}]
    \noindent
    ধরি প্রতিটি গ্যাসের ভর $m$। নাইট্রোজেনের মোল সংখ্যা $n_1 = \frac{m}{28}$ এবং অক্সিজেনের মোল সংখ্যা $n_2 = \frac{m}{32}$। মোট মোল সংখ্যা:
    \[ n_{\text{total}} = n_1 + n_2 = m \left(\frac{1}{28} + \frac{1}{32}\right) = m \left(\frac{8 + 7}{224}\right) = \frac{15 m}{224} \]
    নাইট্রোজেনের মোল ভগ্নাংশ:
    \[ x_{\text{N}_2} = \frac{n_1}{n_{\text{total}}} = \frac{m/28}{15m/224} = \frac{224}{28 \times 15} = \frac{8}{15} \]
    ডাল্টনের আংশিক চাপ সূত্রানুসারে নাইট্রোজেনের আংশিক চাপ:
    \[ P_{\text{N}_2} = x_{\text{N}_2} \times P_{\text{total}} = \frac{8}{15} \times 300\text{ kPa} = \mathbf{160\text{ kPa}} \]
\end{tcolorbox}

\vspace{2em}

% ------------------------------------------
% QUESTION 21 SECTION
% ------------------------------------------
\begin{tcolorbox}[questionbox={প্রশ্ন (Question) 21}]
    \noindent
    স্থির তাপমাত্রায় (isothermal) নির্দিষ্ট ভরের একটি আদর্শ গ্যাসের জন্য $P$ বনাম $1/V$ গ্রাফটি মূলবিন্দুগামী একটি সরলরেখা। এই সরলরেখাটির ঢাল (slope) নিচের কোন রাশির সমান?
    
    \vspace{0.8em}
    \begin{english}
    \noindent\textit{\color{darkgray}
    For a fixed mass of an ideal gas at constant temperature, the plot of $P$ versus $1/V$ is a straight line passing through the origin. What is the slope of this straight line equal to?
    }
    \end{english}
\end{tcolorbox}

\begin{itemize}[label={}, leftmargin=1cm, itemsep=0.5em]
    \item[\textbf{A)}] $n R T$
    \item[\textbf{B)}] $\frac{n R}{T}$
    \item[\textbf{C)}] $\frac{P}{T}$
    \item[\textbf{D)}] $\frac{V}{T}$
\end{itemize}

\vspace{0.5em}
\begin{tcolorbox}[answerbox]
    \textbf{\color{success} উত্তর:} A) $n R T$
\end{tcolorbox}
\vspace{1em}

\begin{tcolorbox}[solutionbox={সমাধান (Solution)}]
    \noindent
    আদর্শ গ্যাস সমীকরণ $P V = n R T$ কে নিম্নরূপে লেখা যায়:
    \[ P = (n R T) \times \left(\frac{1}{V}\right) \]
    এটি $y = m x$ আকারের মূলবিন্দুগামী সরলরেখার সমীকরণ, যেখানে $y = P$ এবং $x = \frac{1}{V}$। সুতরাং, গ্রাফের ঢাল (slope) = $\mathbf{n R T}$।
\end{tcolorbox}

\vspace{2em}

% ------------------------------------------
% QUESTION 22 SECTION
% ------------------------------------------
\begin{tcolorbox}[questionbox={প্রশ্ন (Question) 22}]
    \noindent
    একই তাপমাত্রা ও চাপে একটি অজানা গ্যাস X এর অণুস্রাব হার (rate of effusion) অক্সিজেন গ্যাসের ($M_{\text{O}_2} = 32\text{ g mol}^{-1}$) অণুস্রাব হারের $0.5$ গুণ। গ্যাসটির আণবিক ভর $M_X$ কত এবং একই তাপমাত্রায় অণুর $c_{\text{rms}}$ এর অনুপাত ($\frac{c_{\text{rms},X}}{c_{\text{rms},\text{O}_2}}$) কত?
    
    \vspace{0.8em}
    \begin{english}
    \noindent\textit{\color{darkgray}
    Under identical conditions of temperature and pressure, the rate of effusion of an unknown gas X is $0.5$ times that of Oxygen gas ($M_{\text{O}_2} = 32\text{ g mol}^{-1}$). What is the molar mass of gas X, and what is the ratio of their rms speeds ($c_{\text{rms},X} / c_{\text{rms},\text{O}_2}$) at the same temperature?
    }
    \end{english}
\end{tcolorbox}

\begin{itemize}[label={}, leftmargin=1cm, itemsep=0.5em]
    \item[\textbf{A)}] $M_X = 128\text{ g mol}^{-1}, \frac{c_{\text{rms},X}}{c_{\text{rms},\text{O}_2}} = 0.5$
    \item[\textbf{B)}] $M_X = 64\text{ g mol}^{-1}, \frac{c_{\text{rms},X}}{c_{\text{rms},\text{O}_2}} = 0.707$
    \item[\textbf{C)}] $M_X = 16\text{ g mol}^{-1}, \frac{c_{\text{rms},X}}{c_{\text{rms},\text{O}_2}} = 2.0$
    \item[\textbf{D)}] $M_X = 256\text{ g mol}^{-1}, \frac{c_{\text{rms},X}}{c_{\text{rms},\text{O}_2}} = 0.25$
\end{itemize}

\vspace{0.5em}
\begin{tcolorbox}[answerbox]
    \textbf{\color{success} উত্তর:} A) $M_X = 128\text{ g mol}^{-1}, \frac{c_{\text{rms},X}}{c_{\text{rms},\text{O}_2}} = 0.5$
\end{tcolorbox}
\vspace{1em}

\begin{tcolorbox}[solutionbox={সমাধান (Solution)}]
    \noindent
    গ্রাহামের অণুস্রাব সূত্রানুসারে:
    \[ \frac{r_X}{r_{\text{O}_2}} = \sqrt{\frac{M_{\text{O}_2}}{M_X}} = 0.5 \]
    বর্গ করে পাই:
    \[ \frac{32}{M_X} = 0.25 \implies M_X = \frac{32}{0.25} = \mathbf{128\text{ g mol}^{-1}} \]
    আদর্শ গ্যাসের জন্য $c_{\text{rms}} \propto \frac{1}{\sqrt{M}}$ হওয়ায়:
    \[ \frac{c_{\text{rms}, X}}{c_{\text{rms}, \text{O}_2}} = \sqrt{\frac{M_{\text{O}_2}}{M_X}} = \mathbf{0.5} \]
\end{tcolorbox}

\vspace{2em}

% ------------------------------------------
% QUESTION 23 SECTION
% ------------------------------------------
\begin{tcolorbox}[questionbox={প্রশ্ন (Question) 23}]
    \noindent
    $25^\circ\text{C}$ তাপমাত্রায় বায়ুর আপেক্ষিক আর্দ্রতা $60\%$। যদি বায়ুর তাপমাত্রা কমিয়ে $15^\circ\text{C}$ করা হয়, যা ওই বায়ুর শিশিরাঙ্ক (dew point), তবে $15^\circ\text{C}$ তাপমাত্রায় বায়ুর আপেক্ষিক আর্দ্রতা কত হবে?
    
    \vspace{0.8em}
    \begin{english}
    \noindent\textit{\color{darkgray}
    At $25^\circ\text{C}$, the relative humidity of air is $60\%$. If the air temperature is lowered to $15^\circ\text{C}$, which is the dew point of the air, what will be the relative humidity of the air at $15^\circ\text{C}$?
    }
    \end{english}
\end{tcolorbox}

\begin{itemize}[label={}, leftmargin=1cm, itemsep=0.5em]
    \item[\textbf{A)}] $100\%$
    \item[\textbf{B)}] $60\%$
    \item[\textbf{C)}] $85\%$
    \item[\textbf{D)}] $40\%$
\end{itemize}

\vspace{0.5em}
\begin{tcolorbox}[answerbox]
    \textbf{\color{success} উত্তর:} A) $100\%$
\end{tcolorbox}
\vspace{1em}

\begin{tcolorbox}[solutionbox={সমাধান (Solution)}]
    \noindent
    শিশিরাঙ্ক (Dew Point) হলো সেই তাপমাত্রা যে তাপমাত্রায় বায়ুতে উপস্থিত নির্দিষ্ট পরিমাণ জলীয় বাষ্প দ্বারা বায়ু সম্পৃক্ত (saturated) হয়ে যায়। যেহেতু $15^\circ\text{C}$ হলো এই বায়ুর শিশিরাঙ্ক, তাই তাপমাত্রা হ্রাস পেয়ে $15^\circ\text{C}$ এ পৌঁছালে ওই বাষ্প দ্বারাই বায়ু সম্পূর্ণ সম্পৃক্ত হবে। আর যেকোনো সম্পৃক্ত বায়ুর আপেক্ষিক আর্দ্রতা সর্বদা $\mathbf{100\%}$ হয়।
\end{tcolorbox}

\vspace{2em}

% ------------------------------------------
% QUESTION 73 SECTION
% ------------------------------------------
\begin{tcolorbox}[questionbox={প্রশ্ন (Question) 24}]
    \noindent
    একটি পাত্রে $300\text{ K}$ তাপমাত্রায় এবং $1.0 \times 10^5\text{ Pa}$ চাপে নাইট্রোজেন গ্যাস ($M = 28\text{ g mol}^{-1}, d = 3.7 \times 10^{-10}\text{ m}$) রয়েছে। গ্যাসের তাপমাত্রা অপরিবর্তিত রেখে আয়তন সংকুচিত করে চাপ $2.0 \times 10^5\text{ Pa}$ করা হলে অণুগুলোর দুটি সংঘাতের মধ্যবর্তী গড় সময়কাল (mean collision time, $\tau$) কত হবে? ($k = 1.38 \times 10^{-23}\text{ J K}^{-1}$, $R = 8.314\text{ J mol}^{-1}\text{ K}^{-1}$)
    
    \vspace{0.8em}
    \begin{english}
    \noindent\textit{\color{darkgray}
    A container holds Nitrogen gas ($M = 28\text{ g mol}^{-1}$, $d = 3.7 \times 10^{-10}\text{ m}$) at a temperature of $300\text{ K}$ and a pressure of $1.0 \times 10^5\text{ Pa}$. If the gas is compressed isothermally such that the pressure becomes $2.0 \times 10^5\text{ Pa}$, what will be the mean time between collisions ($\tau$) of the molecules? ($k = 1.38 \times 10^{-23}\text{ J K}^{-1}$, $R = 8.314\text{ J mol}^{-1}\text{ K}^{-1}$)
    }
    \end{english}
\end{tcolorbox}
\begin{itemize}[label={}, leftmargin=1cm, itemsep=0.5em]
    \item[\textbf{A)}] $\tau = 6.58 \times 10^{-11}\text{ s}$
    \item[\textbf{B)}] $\tau = 1.32 \times 10^{-10}\text{ s}$
    \item[\textbf{C)}] $\tau = 3.29 \times 10^{-11}\text{ s}$
    \item[\textbf{D)}] $\tau = 9.87 \times 10^{-11}\text{ s}$
\end{itemize}

\vspace{0.5em}
\begin{tcolorbox}[answerbox]
    \textbf{\color{success} উত্তর:} A) $\tau = 6.58 \times 10^{-11}\text{ s}$
\end{tcolorbox}
\vspace{1em}

\begin{tcolorbox}[solutionbox={সমাধান (Solution)}]
    \noindent
    \textbf{ধাপ ১: মূল গড় বর্গবেগ ($c_{rms}$) নির্ণয়}
    \[ c_{rms} = \sqrt{\frac{3 R T}{M}} = \sqrt{\frac{3 \times 8.314 \times 300}{28 \times 10^{-3}}} = 516.95\text{ m s}^{-1} \]
    
    \vspace{0.4em}
    \noindent\textbf{ধাপ ২: পরিবর্তিত চাপে প্রতি একক আয়তনে অণুর সংখ্যা ($n_V$) নির্ণয়}
    \[ n_V = \frac{P}{k T} = \frac{2.0 \times 10^5}{1.38 \times 10^{-23} \times 300} = 4.8309 \times 10^{25}\text{ m}^{-3} \]
    
    \vspace{0.4em}
    \noindent\textbf{ধাপ ৩: গড় মুক্ত পথ ($\lambda$) ও গড় সময়কাল ($\tau$) হিসাব}
    \[ \lambda = \frac{1}{\sqrt{2} \pi d^2 n_V} = \frac{1}{1.4142 \times 3.1416 \times (3.7 \times 10^{-10})^2 \times (4.8309 \times 10^{25})} = 3.4034 \times 10^{-8}\text{ m} \]
    দুইটি সংঘাতের মধ্যবর্তী গড় সময়কাল:
    \[ \tau = \frac{\lambda}{c_{rms}} = \frac{3.4034 \times 10^{-8}}{516.95} = 6.583 \times 10^{-11}\text{ s} \approx \mathbf{6.58 \times 10^{-11}\text{ s}} \]
\end{tcolorbox}
% ------------------------------------------
% QUESTION 25 SECTION
% ------------------------------------------
\begin{tcolorbox}[questionbox={প্রশ্ন (Question) 25}]
    \noindent
    $27^\circ\text{C}$ তাপমাত্রায় এবং $1.013 \times 10^5\text{ Pa}$ চাপে $1\text{ L}$ পাত্রে থাকা হিলিয়াম গ্যাসের ($M = 4\text{ g mol}^{-1}$) অণুগুলোর মোট অনুবাদ গতিশক্তি (translational kinetic energy) কত? ($R = 8.314\text{ J mol}^{-1}\text{ K}^{-1}$)
    
    \vspace{0.8em}
    \begin{english}
    \noindent\textit{\color{darkgray}
    What is the total translational kinetic energy of Helium gas molecules ($M = 4\text{ g mol}^{-1}$) contained in a $1\text{ L}$ vessel at a temperature of $27^\circ\text{C}$ and a pressure of $1.013 \times 10^5\text{ Pa}$? ($R = 8.314\text{ J mol}^{-1}\text{ K}^{-1}$)
    }
    \end{english}
\end{tcolorbox}

\begin{itemize}[label={}, leftmargin=1cm, itemsep=0.5em]
    \item[\textbf{A)}] $151.95\text{ J}$
    \item[\textbf{B)}] $101.30\text{ J}$
    \item[\textbf{C)}] $253.25\text{ J}$
    \item[\textbf{D)}] $303.90\text{ J}$
\end{itemize}

\vspace{0.5em}
\begin{tcolorbox}[answerbox]
    \textbf{\color{success} উত্তর:} A) $151.95\text{ J}$
\end{tcolorbox}
\vspace{1em}

\begin{tcolorbox}[solutionbox={সমাধান (Solution)}]
    \noindent
    \textbf{ধাপ ১: সূত্র প্রয়োগ} \\
    আদর্শ গ্যাসের অণুগুলোর মোট অনুবাদ গতিশক্তি (Translational Kinetic Energy) নির্ণয়ের সূত্রটি হলো:
    \[ E_k = \frac{3}{2} n R T = \frac{3}{2} P V \]
    
    \vspace{0.4em}
    \noindent\textbf{ধাপ ২: মান প্রতিস্থাপন ও গণনা} \\
    দেওয়া আছে: $P = 1.013 \times 10^5\text{ Pa}$ এবং $V = 1\text{ L} = 1 \times 10^{-3}\text{ m}^3$। মানগুলো সমীকরণে বসিয়ে পাই:
    \begin{align*}
        E_k &= \frac{3}{2} \times (1.013 \times 10^5) \times (1 \times 10^{-3}) \\
        &= 1.5 \times 101.3 = \mathbf{151.95\text{ J}}
    \end{align*}
\end{tcolorbox}
\end{document}